\documentclass[11pt,a4paper]{article}
\usepackage{amsmath,amssymb,amsthm}
\usepackage[margin=2.5cm]{geometry}
\usepackage{hyperref}

\newtheorem{definition}{Definition}[section]
\newtheorem{theorem}[definition]{Theorem}
\newtheorem{proposition}[definition]{Proposition}
\newtheorem{lemma}[definition]{Lemma}
\newtheorem{corollary}[definition]{Corollary}
\newtheorem{conjecture}[definition]{Conjecture}
\newtheorem{remark}[definition]{Remark}

\title{Hyperseed Formalization:\\
  Linguistic Universals as Shadows of Cognitive Structure, Part 3}
\author{ProtomegaTron (automated formalization)}
\date{2026-09-17}

\begin{document}
\maketitle

\begin{abstract}
We formalize the intellectual structure of Ben Goertzel's Substack article
``Linguistic Universals as Shadows of Cognitive Structure, Part 3''
(2026-06-08) within the Hyperseed ontology. The article draws out the
AGI-architecture implications of the cognition-to-language categorical
framework: causal factorization as fiber decomposition, small commutators
as approximate fiber commutativity, semantic frames as transition functions,
graded stability as fiber depth, naturality loss as cocycle condition,
mediator loss as anti-scalar-collapse in coupling, closure loss as
hierarchical cocycle, and linguistic typology as distributed fiber
inspection of cognitive architecture. The concrete training recipe
(category-guided causal-coding Transformers) maps each categorical
requirement onto a differentiable loss with a clean Hyperseed interpretation.
\end{abstract}

\tableofcontents

%% =================================================================
\section{Causal factorization as fiber decomposition}
\label{sec:factorization}
%% =================================================================

\begin{theorem}[Cognitive substrate decomposition --- claims 1--5, 31--32]
\label{thm:factorization}
The article identifies two structural requirements for a language-capable
AGI's cognitive substrate:

\begin{enumerate}
\item \textbf{Causal factorization.} The system's internal state must
  decompose into pieces capturing genuinely separate aspects of the
  world --- different domains, modalities, conceptual neighborhoods,
  skills. Updates targeting one piece should mostly leave others alone.

\item \textbf{Small commutators.} Learning situation A then B should
  produce nearly the same state as B then A, when they touch genuinely
  different parts. If learning Spanish overwrites French, the
  architecture has gotten this wrong.
\end{enumerate}

In Hyperseed terms, causal factorization is \emph{fiber decomposition}:
the cognitive substrate decomposes into fibers over different domains,
and updates are local to the relevant fiber:
\[
  S_{\text{cognitive}} \cong \bigoplus_{d \in \text{Domains}} F_d
  \qquad \text{with} \qquad
  \text{Update}_A \text{ local to } F_{d(A)}
\]

Small commutators are \emph{approximate commutativity of fiber updates}:
updates to independent fibers commute:
\[
  U_A \circ U_B \approx U_B \circ U_A
  \qquad \text{when } F_{d(A)} \cap F_{d(B)} = \emptyset
\]

Non-commutativity (interference between unrelated domains) is a
\emph{fiber defect} --- a failure of the decomposition to be genuine.
The commutator $[U_A, U_B] = U_A \circ U_B - U_B \circ U_A$ measures
the defect; a well-factorized architecture has small commutators for
updates to unrelated domains.

The implementation is deliberately left open: cortical columns, expert
sub-networks, mixture-of-experts, dedicated parameter regions, or
separate symbolic processes. What matters is the relevant separability
--- the fiber decomposition --- not the mechanism that implements it.

Without causal factorization, ``every new thing learned fights every
old thing known.'' This is catastrophic forgetting viewed as a fiber
defect: the fibers are not genuinely separate, so updates to one
fiber corrupt others. The factorization theorem behind broad universals
\emph{requires} this property --- it's the structural precondition for
cross-domain generalization.
\end{theorem}

%% =================================================================
\section{Semantic frames as transition functions}
\label{sec:frames}
%% =================================================================

\begin{definition}[The cognition-to-language bundle --- claims 6--8, 33]
\label{def:frames}
The article identifies semantic frames (event frames, force-dynamic
frames, argument-role frames, person and number frames, possession
frames, attention frames) as the intermediate layer bridging cognition
and language. A single causal module typically supports a cluster of
related conceptual schemas; a single linguistic feature externalizes
only a piece of one schema.

In Hyperseed terms, semantic frames are the \emph{transition functions}
of the cognition-to-language bundle:
\[
  E_{\text{cog}} \xrightarrow{g_{\text{frame}}} E_{\text{lang}}
\]

The frame layer is what makes the bundle well-defined: without it,
the cognition-to-language map is ``too rigid to do useful work'' ---
in bundle terms, the transition functions would have to be the identity
(every cognitive structure maps directly to a linguistic structure),
which is empirically false. The frame layer provides the necessary
degrees of freedom for the transition functions to vary across
contexts while preserving the bundle structure.

The frame layer can be ``discrete, continuous, hybrid, or
neural-symbolic'' --- what matters is that frame transformations and
linguistic transformations be comparable through a loss. This is the
implementation-agnostic requirement: the transition functions must
exist as a well-defined mathematical structure, but their realization
is unconstrained.
\end{definition}

%% =================================================================
\section{Sparse routing and anti-scalar-collapse in coupling}
\label{sec:sparse}
%% =================================================================

\begin{proposition}[Explicit mediators --- claims 9--11, 23, 36]
\label{prop:sparse}
The article requires sparse externalization routing: connections
between cognitive and linguistic structure should be sparse --- a
small number of dedicated high-traffic pathways rather than dense
coupling everywhere. The evidence: the small number of broad
cross-module universals that survive genealogical and areal
correction (case marking, agreement morphology, argument-structure
resolution) are precisely the high-traffic mediators.

The mediator loss formalizes this: two distant module families should
be conditionally independent given an explicit mediator module. The
architecture is forbidden from implementing broad dependencies as
diffuse dense coupling between every relevant head; a small number of
explicit mediator heads must absorb that traffic.

In Hyperseed terms, this is \emph{anti-scalar-collapse applied to
inter-module coupling}: the coupling structure must be explicit and
inspectable, not collapsed into diffuse dense connections. Dense
coupling is scalar collapse of inter-module structure --- it hides
the dependency structure inside a matrix of small interactions
rather than exposing it as a small number of identifiable pathways.

The cost of dense coupling: interference, inference cost, and
catastrophic forgetting. These are the practical consequences of
scalar collapse in coupling --- when everything couples to everything,
there's no way to update one part without affecting all others, and
there's no way to inspect the dependency structure to understand why
a change propagated.
\end{proposition}

%% =================================================================
\section{Graded stability as fiber depth}
\label{sec:stability}
%% =================================================================

\begin{definition}[Three layers of universality --- claims 12, 15, 25, 34]
\label{def:stability}
The article identifies three layers of linguistic universality,
corresponding to three levels of cognitive stability:

\begin{enumerate}
\item \textbf{Hierarchical universals (bedrock).} Person, number,
  case, and accessibility hierarchies. These name cognitive bedrock
  --- the most stable, most protected, most deeply encoded structures.
  Slow learning rates, strong consolidation.

\item \textbf{Narrow word-order regularities (consolidated preferences).}
  These name a layer of externalization preference that is stable but
  not as deeply protected as the hierarchies. Moderate stability.

\item \textbf{Long tail of context-bound universals (outer shell).}
  Where most actual cognitive variation lives. High plasticity,
  fast adaptation, context-sensitivity.
\end{enumerate}

In Hyperseed terms, this is \emph{fiber depth}: core universals live
in deep (protected) fiber layers; context-local patterns live in
shallow (plastic) layers. The stability schedule --- slow learning
rates for deep layers, fast learning rates for shallow layers ---
mechanically distinguishes the rigid inner kernel from the flexible
outer shells by what kind of universal structure they represent.

This is the same graded-fiber-depth structure as the four-layer
architecture in the AGI Catastrophe Part 2 article, applied to
cognitive architecture rather than security architecture: the
deepest layers are most protected, and the system accesses deeper
layers only when the task requires it.
\end{definition}

%% =================================================================
\section{Training losses as categorical conditions}
\label{sec:losses}
%% =================================================================

\begin{theorem}[Category-guided training recipe --- claims 17--27, 35, 37]
\label{thm:losses}
The article's engineering contribution: a concrete training recipe
for Transformer neural networks in which each loss has a clean
categorical (and therefore Hyperseed) interpretation.

\begin{enumerate}
\item \textbf{Naturality loss $\to$ cocycle condition.} Updating the
  hidden state then reading out linguistic structure should give the
  same result as reading out then applying the corresponding
  linguistic update:
  \[
    \text{Readout} \circ \text{Update}_{\text{cog}} \approx
    \text{Update}_{\text{lang}} \circ \text{Readout}
  \]
  This is the cocycle condition on the cognition-to-language bundle:
  the transition functions must compose consistently. If the diagram
  doesn't commute, the cognition-to-language map is not a well-defined
  bundle morphism --- the linguistic readout is not faithfully tracking
  the cognitive state.

\item \textbf{Closure loss $\to$ hierarchical cocycle.} Detecting a
  marked feature (dual number, third-person object agreement) implies
  detecting every feature below it on the hierarchy. This is a
  hierarchical cocycle condition: the fiber structure at each level
  must be consistent with all levels below:
  \[
    \text{Detect}(f_{\text{marked}}) \implies
    \bigwedge_{f' \leq f_{\text{marked}}} \text{Detect}(f')
  \]

\item \textbf{Mediator loss $\to$ anti-scalar-collapse.} Two distant
  module families should be conditionally independent given an
  explicit mediator. The dependency structure must be explicit and
  inspectable.

\item \textbf{Low-frustration loss $\to$ sparse signed graph.}
  Linearization choices organize into a sparse signed graph with
  mostly compatible local pulls. This reflects the constraint that
  externalization preferences should be locally consistent even if
  globally they can't all be simultaneously satisfied.

\item \textbf{Stability schedule $\to$ fiber depth.} Slow learning
  for deep (high-stability) layers, fast learning for shallow
  (context-local) layers. The depth assignment is not hand-coded
  but learned from which structures the other losses identify as
  stable.
\end{enumerate}

The key design principle: ``linguistic universals are not curiosities
but training targets --- a route by which thousands of years of
cross-cultural averaging done by human language itself can be fed
back into machine learning systems as architectural and training-time
signals about what the latent causal structure of cognition should
look like.''
\end{theorem}

%% =================================================================
\section{Linguistic typology as distributed fiber inspection}
\label{sec:typology}
%% =================================================================

\begin{proposition}[Cross-cultural averaging --- claims 13--14, 16, 38]
\label{prop:typology}
The article makes a strong methodological claim: linguistic typology,
once corrected for genealogy and area, gives cleaner evidence about
what's stable in the human cognitive machinery than most direct
cognitive experiments can.

In Hyperseed terms, this is \emph{distributed fiber inspection}:
many independent languages, developed over thousands of years by
millions of speakers across diverse cultures, are all inspecting
the same underlying cognitive fiber through their surface
externalizations. The typological record is the result of this
distributed inspection --- features that appear across unrelated
languages in unrelated areas are features of the shared fiber,
not accidents of particular sections.

The power of this approach: the cross-cultural averaging that
``would normally take heroic effort for any individual (human or AI)
mind to do has already been done by language itself.'' Each language
is a section of the cognition-to-language bundle; the typological
universals are the fiber properties that persist across all sections.

This connects to the broader Hyperseed theme of provenance through
diversity: the more independent the inspectors (languages) and the
stronger the correction for shared ancestry (phylogenetic control),
the higher the confidence that a detected universal is a genuine
fiber property rather than a section artifact.
\end{proposition}

%% =================================================================
\section{Convergence of independent lines}
\label{sec:convergence}
%% =================================================================

\begin{theorem}[Structural identity --- claims 28--30]
\label{thm:convergence}
The article's strongest claim: two independent lines of work ---
linguistic typology (empirical) and brain-like continual learning
(theoretical) --- converge on the same mathematical structure:

\begin{itemize}
\item Closure systems with diachronic Lyapunov flow.
\item Factorization theorems with mediator exceptions.
\item Sparse signed interaction graphs.
\item Graded stability under transition dynamics.
\item Gluing of context-indexed local closure operators.
\end{itemize}

This convergence is evidence that both lines are tracking the same
underlying phenomenon: the mathematical structure of cognition itself,
as it manifests both in the typological record (downstream linguistic
consequences) and in the design requirements for artificial cognitive
systems (upstream architectural constraints).

The Occam-like principle that picks out the right cognitive explanation
of a linguistic pattern is formalized as an algebraic order on
explanations combining five criteria: indistinction, fit, cost,
confusion, and double-counting. The result that causal explanations
beat correlation-only explanations ``falls out as a theorem in that
order'' --- it's not a preference but a consequence of the algebraic
structure.

In Hyperseed terms, this convergence is evidence for the Hyperseed
ontology itself: if two independent investigations of cognitive
architecture arrive at the same fiber-bundle structure (decomposition,
commutativity, transition functions, depth, cocycle conditions), that
structure is a genuine property of cognition, not an artifact of
either investigation's methodology.
\end{theorem}

%% =================================================================
\section{Cross-references}
%% =================================================================

\begin{remark}[Connections to other formalizations]
\begin{itemize}
\item \textbf{Linguistic Universals Parts 1--2}: companion articles
  developing the categorical framework and its empirical grounding.
\item \textbf{Avoiding AGI Catastrophe Part 2} (2026-06-10): graded
  fiber depth. The four-layer security architecture has the same
  graded-depth structure as the three-layer stability architecture
  here.
\item \textbf{Seeding RSI} (2026-07-28): Hyperon architecture. The
  causal-coding Transformer recipe is a complement to the Hyperon
  neural-symbolic architecture --- both aim at language-capable AGI
  with inspectable internal structure.
\item \textbf{The Orchard Bug} (2026-06-05): cocycle defects. The
  naturality loss is designed to prevent the cognitive analog of
  the Orchard bug --- composition failures where individual modules
  work but their interaction doesn't.
\item \textbf{Un-Jailbreakable AI} (2026-07-10): generate-and-verify.
  The closure loss and mediator loss implement a form of
  generate-and-verify: the model generates internal representations
  and the losses verify that they satisfy the categorical constraints.
\item \textbf{What Should an AGI Curriculum Look Like?} (2026-07-22):
  learning architecture. The graded stability schedule connects to
  the curriculum design question --- what should be learned first
  (deep, stable structures) vs.\ last (shallow, context-local ones).
\item \textbf{$(f,c)$-lossy} (LTM 5a4d150b): anti-scalar-collapse.
  The mediator loss is anti-scalar-collapse applied to inter-module
  coupling --- the same principle applied to a different substrate.
\item \textbf{Non-invertible 2-cells} (LTM 85c02e32): composition
  structure. The naturality loss ensures that cognitive-to-linguistic
  composition is well-defined --- preventing the non-invertible
  witness from infecting the mapping.
\end{itemize}
\end{remark}

\begin{conjecture}[Universality of the fiber-depth pattern]
Conjecture: the three-layer stability pattern (bedrock / consolidated
preferences / context-local residue) is not specific to language but
appears in every cognitive domain with sufficient developmental depth.
Visual perception (basic features / object categories / scene
interpretation), motor control (reflexes / skills / strategies),
and social cognition (basic emotion recognition / social norms /
cultural conventions) should all exhibit the same graded-stability
structure with the same mathematical properties (closure at the
deepest level, factorization at the middle level, sparse signed
interaction at the surface level).

If this conjecture holds, the linguistic universals training recipe
generalizes: any domain's ``universals'' (stable cross-context
features) can serve as training targets for the corresponding
cognitive module, using the same categorical loss functions with
domain-appropriate probes.
\end{conjecture}

\end{document}
